% ============================================================
% conclusion.tex — Kết luận
% ============================================================

\chapter*{KẾT LUẬN}
\addcontentsline{toc}{chapter}{KẾT LUẬN}

% ---------------------------------------------------------------
\section*{1. Tóm tắt}
\addcontentsline{toc}{section}{1. Tóm tắt}
% ---------------------------------------------------------------

Đề cương nghiên cứu đã trình bày bài toán, phương pháp và kế hoạch thực hiện cho đề tài \textit{``\detai''}. Đề tài hướng đến xây dựng một prototype nghiên cứu end-to-end cho bài toán dự đoán lan truyền khói cấp phòng/khu vực trong tòa nhà cao tầng, tích hợp các thành phần:

\begin{itemize}
  \item Module BIM/IFC-to-Graph chuyển đổi tự động mô hình tòa nhà thành đồ thị.
  \item Mô phỏng batch bằng CFAST kết hợp kiểm chứng FDS.
  \item Hệ thống IoT giả lập 3 lớp tạo dữ liệu cảm biến thực tế.
  \item Mô hình AI surrogate dạng Temporal GNN kết hợp thông tin đồ thị, chuỗi thời gian và điều kiện biên.
  \item Dashboard trực quan hóa và replay kết quả.
\end{itemize}

% ---------------------------------------------------------------
\section*{2. Tính mới và đóng góp dự kiến}
\addcontentsline{toc}{section}{2. Tính mới và đóng góp dự kiến}
% ---------------------------------------------------------------

Tính mới của đề tài nằm ở \textbf{pipeline tích hợp}, không phải ở việc phát minh riêng lẻ từng thành phần. Cụ thể:

\begin{enumerate}
  \item \textbf{Pipeline BIM → Simulation → AI:} Kết nối liền mạch từ mô hình BIM đến dự đoán AI, cho phép tự động hóa toàn bộ quy trình.
  \item \textbf{Building graph từ BIM/IFC:} Tận dụng thông tin kiến trúc có sẵn thay vì xây dựng mô hình thủ công.
  \item \textbf{Virtual IoT 3 lớp:} Tạo dữ liệu huấn luyện AI phản ánh tính bất hoàn hảo của cảm biến thực tế.
  \item \textbf{Boundary-Aware Temporal GNN:} Mô hình AI tích hợp đồng thời cấu trúc đồ thị, chuỗi thời gian, điều kiện biên và trạng thái hệ thống PCCC.
  \item \textbf{Đánh giá đa chiều:} So sánh 7 mức mô hình, nhiều chế độ chia dữ liệu, ablation study, và kiểm tra robustness.
\end{enumerate}

% ---------------------------------------------------------------
\section*{3. Giới hạn nghiên cứu}
\addcontentsline{toc}{section}{3. Giới hạn nghiên cứu}
% ---------------------------------------------------------------

Nghiên cứu thừa nhận các giới hạn sau:

\begin{table}[H]
\centering
\begin{tabularx}{\textwidth}{l X}
\toprule
\textbf{Giới hạn} & \textbf{Ý nghĩa} \\
\midrule
Không có dữ liệu cháy thật      & Đánh giá dựa trên output mô phỏng \\
Không có IoT phần cứng           & IoT là mô phỏng hoàn toàn \\
CFAST là zone-level              & Không dự đoán trường 3D \\
FDS chỉ kiểm chứng chọn lọc     & Không phải nguồn dataset chính \\
Nhãn hazard có thể là proxy      & Trừ khi có ngưỡng chuẩn \\
AI là surrogate                  & Không thay thế CFD \\
Dashboard là demo nghiên cứu     & Không phải hệ thống điều khiển PCCC \\
Không có kết luận pháp lý        & Không tuyên bố đạt thẩm duyệt \\
\bottomrule
\end{tabularx}
\end{table}

% ---------------------------------------------------------------
\section*{4. Hướng phát triển tương lai}
\addcontentsline{toc}{section}{4. Hướng phát triển tương lai}
% ---------------------------------------------------------------

\begin{itemize}
  \item Mở rộng FDS validation (8--12 case hoặc toàn tòa nếu đủ phần cứng).
  \item Tích hợp IoT phần cứng thật (cảm biến khói, nhiệt, CO).
  \item Mô hình hóa chi tiết hệ thống HVAC (duct, damper).
  \item Tích hợp mô hình thoát nạn (evacuation) để tính ASET/RSET.
  \item Mô phỏng CFD chi tiết cho atrium và không gian lớn.
  \item Phát triển hệ thống điều khiển PCCC thông minh dựa trên dự đoán AI.
  \item Kiểm chứng trên nhiều tòa nhà thực tế để đánh giá tính tổng quát.
\end{itemize}
